% 选题阶段工作笔记；正式论文确定后再迁移至 APS REVTeX 模板。
\documentclass[UTF8,11pt,a4paper]{ctexart}
\usepackage[margin=2.6cm]{geometry}
\usepackage{amsmath,amssymb}
\usepackage[colorlinks=true,linkcolor=blue,citecolor=blue,urlcolor=blue]{hyperref}
\usepackage[backend=biber,style=numeric,sorting=none,doi=true,url=false,eprint=true]{biblatex}
\renewcommand*{\bibfont}{\small}
\setlength{\bibitemsep}{0pt}
\addbibresource{../../literature/references.bib}

\title{相互作用 Dicke--Ising 模型有限温度临界结构\\[0.4em]
\large 候选选题、淘汰边界及国内外研究现状（草稿 v0.3）}
\author{Quantum Optics Thesis 工作笔记}
\date{2026 年 11 月 17 日}

\begin{document}
\maketitle

\begin{abstract}
本文是毕业设计准备阶段的文献与模型笔记，尚非已确定的论文题目或研究结论。
项目以形成投稿潜力为选题筛选目标；当前首要候选是与近期零温 Dicke--Ising 工作轴向一致的
一级/二级边界和三临界结构如何延伸到有限温度，但该文献空白仍待系统查重和方法可行性验证。
静态单模、同一正交分量上的 $A^2$ 与反压缩项可由精确压缩变换吸收，
故不再把“反压缩与自旋相互作用的简单叠加”作为新机制；它只保留为解析基线、受约束控制路径和原光子观测量问题。
对外部泵浦产生反压缩的装置，实验室系通常属于受驱动体系，应与真正无泵浦平衡态区分。
\end{abstract}

\section{问题边界与基本模型}

以单模光场和 $N$ 个自旋为起点，取 $\hbar=1$、$s_i^\alpha=\sigma_i^\alpha/2$，写成
\begin{equation}
 H=\omega a^\dagger a+H_{\mathrm{spin}}(\Delta,J)
 +\frac{2g}{\sqrt N}(a+a^\dagger)S_x
 +[D(g)-\xi](a+a^\dagger)^2,
 \qquad S_x=\sum_{i=1}^{N}s_i^x .
 \label{eq:model}
\end{equation}
其中 $D(g)$ 是由选定微观平台决定的 $A^2$ 系数，$\xi$ 是暂时假设可独立控制的反压缩强度。
目前优先选最近邻 Ising 链作为 $H_{\mathrm{spin}}$；自旋相互作用的符号与归一化在开始计算前统一固定。
模型~\eqref{eq:model} 存在正常光场基态至少要求
\begin{equation}
  \omega+4[D(g)-\xi]>0 . \label{eq:stability}
\end{equation}
满足这一条件时，光子压缩变换给出
\begin{equation}
 \omega_s=\sqrt{\omega\{\omega+4[D(g)-\xi]\}},\qquad
 g_s=g\left(\frac{\omega}{\omega+4[D(g)-\xi]}\right)^{1/4}.
 \label{eq:squeezed}
\end{equation}
因此静态单模的二次反压缩项不会单凭自身生成一种新的哈密顿量类别。
受约束的 $(g,J,\xi)$ 路径及原光子观测量可在主问题通过查重后作为从属分析，
不能独立承担投稿新颖性。

\section{国内及国内参与的相关研究}

Zhang 等研究超导量子比特阵列与腔模构成的 Dicke--Ising 模型，
预测正常、超辐射、反铁磁正常及反铁磁超辐射共存等相，说明磁序与光场序可同时存在\cite{Zhang2014CircuitIsing}。
这为研究自旋相互作用提供早期理论背景，也提醒我们不能把共存相直接宣称为新发现。

Zhu、Ping、Yang 和 Agarwal 讨论参量放大产生的压缩光对集体超辐射相变与三临界点的控制\cite{Zhu2020SqueezedLight}。
这里的外部泵浦是判断物理体系是否真正处于热平衡的重要条件。
Chen 等使用核磁共振量子模拟器制备含 $A^2$ 与反压缩项的有效 Rabi 模型基态，
报告反向超辐射转变、零点涨落增强及有限频率标度\cite{Chen2021Antisqueezing}。
这项实验研究的是受控模拟器的有效哈密顿量；在引言和讨论中需与真实、无驱动腔体系的热平衡相变分开表述。

Zhu 等随后在双传播方向腔模和定向参量泵浦的模型中讨论非互易、不同阶数的超辐射稳态相变\cite{Zhu2024Nonreciprocal}。
该工作展示开放系统路线的可能性，但其外部泵浦、损耗及模间耦合使它不能直接作为式~\eqref{eq:model} 的平衡态先例。

\section{国际研究现状}

对普通电偶极腔 QED，Nataf 与 Ciuti 依据振子强度求和规则分析时间独立哈密顿量的超辐射禁戒条件，
同时指出某些电路 QED 实现的约束不同\cite{Nataf2010NoGo}。
因此任何关于“绕开禁戒”的结论都必须说明具体平台、规范约束与反压缩的物理来源。
Peng 等系统讨论 Dicke 类模型在零温和有限温度下的超辐射相变，说明热平衡研究并不局限于基态\cite{Peng2019Unified}。
Gammelmark 与 Mølmer 早在 2011 年便对量子 Ising 链和单模腔的有限温度自由能及一、二阶超辐射相界开展研究；
其自旋轴的相对方向与式~\eqref{eq:model} 拟用模型不同，需逐式核对\cite{Gammelmark2011Thermal}。
Otake 与 Bamba 近期对腔耦合一维经典 Ising 链给出了有限温度超辐射相变和临界温度的精确结果；
该模型并未包含本文拟研究的量子自旋涨落与独立反压缩项，但构成直接的有限温度先例\cite{Otake2026FiniteT}。
ErFeO$_3$ 中也已有磁振子型热平衡超辐射转变的材料理论与光谱证据；
该自旋交换平台不适用普通电偶极腔光子的 $A^2$ 问题\cite{Bamba2022Magnonic,Kim2025Magnonic}。
与闭合平衡问题相对，Baumann 等在外部泵浦的开放腔中实验观察 Dicke 型相变，
是比较受驱动路线时的重要实验背景\cite{Baumann2010OpenDicke}。

在相互作用自旋方向，Rohn 等分析一维光诱导量子横场 Ising 链的磁序改变与直接相变\cite{Rohn2020QuantizedIsing}；
Schellenberger 与 Schmidt 给出关联光--物质系统正常相低能部分的映射方法\cite{Schellenberger2024Mapping}。
Mendonça、Jachymski 与 Wang 研究 Dicke--Ising/XXZ 的基态，报告自旋相互作用对相界及
XXZ 中间区域次线性光子占据的影响\cite{Mendonca2025Matter}。
Langheld、Hörmann 与 Schmidt 则用大尺度量子蒙特卡洛给出 Dicke--Ising 的定量相图，
包括较窄的反铁磁与超辐射共存区；其补充材料第 IV 节已将 $A^2$ 项通过压缩变换吸收，
并讨论特定求和规则约束下的可达区域\cite{Langheld2025Wormhole}。
Hörmann 等针对前一组的 Ising 相图提出公开质疑，指出共存区和部分跃迁阶数的遗漏\cite{Hormann2025Comment}。
这项争议要求后续计算对照大系统数据与解析极限，而非只引用一篇论文的图。
Leibig 等在已发表于 \emph{SciPost Physics} 的工作中，进一步用热力学极限的自洽纯自旋模型和 NLCE+DMRG 计算零温相图，
确认狭窄的反铁磁--超辐射共存区；这也是本课题数值方法的直接基准\cite{Leibig2026SelfConsistent}。
Hörmann 的 2026 年预印本进一步报告 Dicke--Ising 模型的零温解析相界、三临界条件及反铁磁四相交会附近的结构，
并将有限温度热相界延伸列为尚未完成的方向\cite{Hormann2026Folds}。
这些属于文献作者的结果与展望，并非本项目已复现的计算；仅重绘零温相图不能作为独立创新点。

更近的研究显示，直接耦合物质临界模可降低超辐射阈值并增强本征压缩\cite{Sur2026Critical}。
因此“自旋相互作用与光场压缩协同”不能作为笼统的新颖性主张；在当前项目中，该表述已退出活动核心主张。

\section{当前候选问题与停走边界}

候选主问题是在热力学极限下，以配分函数 $Z=\operatorname{Tr}e^{-\beta H}$ 和自由能
$F=-\beta^{-1}\ln Z$ 研究与近期零温模型相同的 $zz$ 交换、$z$ 纵场和 $x$ 腔耦合体系，
追踪零温一级/二级边界及三临界点向 $T>0$ 的延伸。一级相界必须由所有候选分支的全局自由能比较确定；
有限系统热精确对角化只用于极限和程序检查，不能把平滑交叉称为热相变。

当前定向核查只说明该问题值得完成文献缺口审计，并未证明首次性。
Gammelmark--Mølmer 已给出另一自旋轴向的有限温一、二阶超辐射边界，
Langheld、Leibig、Koziol 与 Hörmann 已显著收窄零温三临界问题。
若找到相同哈密顿量的完整有限温相图，或现有资源无法控制热力学极限与全局分支，本候选应停止或降级。

只有主问题通过上述门槛后，才加入 $D(g)$、$\xi$ 受约束路径，并计算原光场占据
$\langle a^\dagger a\rangle/N$、正交分量方差及均匀/交错磁化。
单模二次项的压缩变换在任意温度保持能谱和配分函数等价，不能把重参数化包装成新相。

\section{当前待办与核查点}

\begin{enumerate}
  \item 完成 Gammelmark--Mølmer、Langheld、Leibig、Koziol 与 Hörmann 的逐式差异矩阵及双向引文检索；不使用“首次”措辞。
  \item 比较有限温自洽自由能、热精确对角化及至少一种可扩展热力学方法的资源和误差，形成方法可行性备忘。
  \item 与导师确认题目、时间表和可用计算资源；上述两项未通过前不启动大规模参数扫描。
  \item 统一 $J$ 的符号、边界条件、光场截断和自旋归一化；复现已知零温铁磁与反铁磁截面。
  \item 继续追踪 2025--2026 年 Ising 相图、有限温与耗散工作的后续版本，更新主张台账和引用库。
\end{enumerate}

\printbibliography[title={参考文献}]
\end{document}
